\begin{vidu}
	Một ấm đun nước có công suất 500 W chứa 300 g nước ở $20^\circ$.
	\begin{itemize}
		\item[a)] Tính thời gian cần thiết để đun nước trong ấm đạt đến nhiệt độ sôi.
		\item[b)] Sau khi nước đến nhiệt độ sôi, người ta để ấm tiếp tục đun nước sôi trong 2 phút. Tính khối lượng nước còn lại trong ấm.
	\end{itemize}
	\loigiai{
		\begin{itemize}
			\item 
			\item 
			\item 
			\item 
			\item 
			\item 
		\end{itemize}
	}
\end{vidu}
\begin{vidu}
	Một nhà máy thép mỗi lần luyện được 35 tấn thép. Cho nhiệt nóng chảy riêng của thép là $2,77.10^5\ J/kg$. 
	\begin{itemize}
		\item[a)] Tính nhiệt lượng cần cung cấp để làm nóng chảy thép trong mỗi lần luyện của nhà máy ở nhiệt độ nóng chảy.
		\item[b)] Giả sử nhà máy sử dụng khí đốt để nấu chảy thép trong lò thổi (nồi nấu thép). Biết khi đốt cháy hoàn toàn 1 kg khí đốt thì nhiệt lượng toả ra là $44.10^6\ J$. Xác định lượng khí đốt cần sử dụng.
	\end{itemize}	
	\loigiai{
		\begin{itemize}
			\item 
			\item 
			\item 
			\item 
			\item 
			\item 
		\end{itemize}
	}
\end{vidu}
\begin{vidu}
	\begin{itemize}
		\item[a)] Một ấm điện công suất 1000 W. Tính gần đúng thời gian cần thiết để đun 300 g nước có nhiệt độ ban đầu là $20^\circ C$ đến khi sôi ở áp suất tiêu chuẩn. Tại sao kết quả chỉ được coi là gần đúng?
		\item[b)] Nếu để nước trong ấm sôi thêm 2 phút thì lượng nước còn lại trong ấm là bao nhiêu? Lấy nhiệt dung riêng và nhiệt hoá hơi riêng của nước là $c=4180.10^3\ J/kg.K$ và $L=2,26.10^6\ J/kg$.
	\end{itemize}
	\loigiai{
		\begin{itemize}
			\item 
			\item 
			\item 
			\item 
			\item 
			\item 
		\end{itemize}
	}
\end{vidu}

\loai{Tính nhiệt nóng chảy}
\begin{vidu}%[1P7-4.2G][Loai1]
	Tính nhiệt lượng Q cần cung cấp để làm nóng chảy 500 g nước đá ở $0^\circ C$. Biết nhiệt nóng chảy riêng của nước đá là $3,4.10^5\ J/kg$.
	\choice
	{$Q=7.10^7\ J$}
	{\True $Q=17.10^4\ J$}
	{$Q=17.10^5\ J$}
	{$Q=17.10^6\ J$}
	\loigiai{
		Ta có: $Q=m.\lambda =0,5.3,4.10^5=17.10^4\ J$.
		\begin{note}
			Vì nhiệt độ $0^\circ C$ là nhiệt độ nóng chảy của nước đá rồi nên nhiệt lượng Q cấp vào sẽ chỉ dùng để làm tan nước đá mà không làm thay đổi nhiệt độ của khối nước đá.	
		\end{note}
	}
\end{vidu}

\loai{Tính nhiệt lượng cung cấp để làm nóng chảy}
\begin{vidu}%[1P7-4.2G][Loai2]
	Một thỏi nhôm có khối lượng 1,0 kg ở $8^\circ C$. Tính nhiệt lượng Q cần cung cấp để làm nóng chảy hoàn toàn thỏi nhôm này. Nhôm nóng chảy ở $658^\circ C$, nhiệt nóng chảy riêng của nhôm là $3,9.10^5\ J/Kg$ và nhiệt dung riêng của nhôm là $880\ J/kg.K$
	\choice
	{$Q\approx 5,9.10^6\ J$}
	{$Q\approx 59.10^4\ J$}
	{$Q\approx 4,47.10^5\ J$}
	{\True $Q\approx 9,62.10^5\ J$}
	\loigiai{
		Ta có: $Q=mc.\Delta t+m.\lambda \Rightarrow Q=1.880\left({658-8}\right)+1.3,9.10^5=9,62.10^5\ J$.
	}
\end{vidu}

\loai{Tính nhiệt lượng cần cung cấp để làm hoá hơi}
\begin{vidu}
	Bạn Tùng đun một ấm nước đầy có dung tích 1,5 lít bằng bếp gas. Do sơ suất nên bạn quên không tắt bếp khi nước sôi. Tính nhiệt lượng cần thiết để làm hoá hơi hoàn toàn nước trong ấm. Biết nhiệt hoá hơi riêng của nước là $2,3.10^6\ J/kg$.
	\loigiai{
		\begin{itemize}
			\item 
			\item 
			\item 
			\item 
			\item 
			\item 
		\end{itemize}
	}
\end{vidu}
